\chapter{Bị Từ Chối, Bị Chê, Bị Thay Thế}
\markboth{Bị Từ Chối, Bị Chê, Bị Thay Thế}{Rejected, Criticized, Replaced}

\section{Không Được Chọn Là Chuyện Rất Bình Thường}
\begin{truyen}{Không Được Chọn Là Chuyện Rất Bình Thường}{Not Being Chosen Is Very Normal}
\chuhoa{N}{hiều người trẻ sốc mạnh trong lần đầu bị từ chối xin việc, xin học bổng, hay xin một cơ hội tốt.}
\textit{(Many young people are deeply shocked the first time they are rejected for a job, a scholarship, or a good chance.)}

Họ tưởng mình làm chưa đủ.
\textit{(They think they did not do enough.)}

Đôi khi đúng.
\textit{(Sometimes that is true.)}

Nhưng cũng nhiều khi chỉ là hôm đó người ta chọn người khác hợp hơn.
\textit{(But often someone else simply fit better that day.)}

Không được chọn đau vì nó chạm vào cái tôi.
\textit{(Not being chosen hurts because it touches the ego.)}

Nhưng nếu học được cách chịu đựng nó, em sẽ bớt gãy hơn rất nhiều trong đời.
\textit{(But if you learn how to bear it, life will break you much less.)}

\end{truyen}

\begin{baihoc}
Bị từ chối không nói hết giá trị của em. Nó chỉ nói rằng cơ hội đó không thuộc về em vào lúc đó.
\textit{(Rejection does not tell the whole truth about your worth. It only says that this chance was not yours at that time.)}
\end{baihoc}

\begin{ghinhoanh}
\textbf{Short English to remember:}

Rejection does not tell the whole truth about your worth.

It only says that this chance was not yours at that time.

\end{ghinhoanh}

\begin{tuhocnhanh}
\textbf{Cau hoi tu hoc / Self-study questions}

1. Van de chinh la gi? \textit{What is the main problem?}

2. Bai hoc thuc te la gi? \textit{What is the practical lesson?}

3. Mot cau tieng Anh em muon nho la cau nao? \textit{Which English sentence do you want to remember?}
\end{tuhocnhanh}
\ngancach

\section{Bị Chê Không Luôn Là Bị Xúc Phạm}
\begin{truyen}{Bị Chê Không Luôn Là Bị Xúc Phạm}{Criticism Is Not Always an Insult}
\chuhoa{Đ}{i học lâu làm nhiều người quen với việc lỗi sai vẫn được nương.}
\textit{(Long years in school make many people used to having mistakes softened.)}

Ra đời rồi, lời góp ý có thể ngắn, thẳng, và khó nghe hơn.
\textit{(In adult life, feedback can be shorter, sharper, and harder to hear.)}

Nếu cái tôi quá lớn, em sẽ thấy mọi lời chê đều là đòn đánh vào mình.
\textit{(If the ego is too large, every criticism feels like a personal attack.)}

Nhưng bớt tự ái một chút, em sẽ thấy nhiều lời khó nghe chính là thứ giúp mình lớn lên nhanh nhất.
\textit{(But if you become a little less defensive, you will see that many hard comments are what help you grow fastest.)}

\end{truyen}

\begin{baihoc}
Không phải lời chê nào cũng đúng. Nhưng người trưởng thành biết lọc phần có ích thay vì chỉ nổi nóng.
\textit{(Not every criticism is correct. But a mature person knows how to keep the useful part instead of only getting angry.)}
\end{baihoc}

\begin{ghinhoanh}
\textbf{Short English to remember:}

Not every criticism is correct.

But a mature person knows how to keep the useful part instead of only getting angry.

\end{ghinhoanh}

\begin{tuhocnhanh}
\textbf{Cau hoi tu hoc / Self-study questions}

1. Van de chinh la gi? \textit{What is the main problem?}

2. Bai hoc thuc te la gi? \textit{What is the practical lesson?}

3. Mot cau tieng Anh em muon nho la cau nao? \textit{Which English sentence do you want to remember?}
\end{tuhocnhanh}
\ngancach

\section{Bị Thay Thế Là Một Bài Học Tỉnh Người}
\begin{truyen}{Bị Thay Thế Là Một Bài Học Tỉnh Người}{Being Replaced Is a Waking Lesson}
\chuhoa{C}{ó người trẻ tưởng chỉ cần mình từng làm tốt là vị trí sẽ luôn là của mình.}
\textit{(Some young people think that if they once did well, the position will always remain theirs.)}

Nhưng đời thay người rất nhanh khi em ngừng học, ngừng giữ chất lượng, hay ngừng tiến lên.
\textit{(But life replaces people quickly when they stop learning, stop keeping quality, or stop moving forward.)}

Bị thay thế nghe đau, nhưng nó nhắc em rằng không có chỗ đứng nào nên được ngủ quên trên đó.
\textit{(Being replaced sounds painful, but it reminds you that no position should become a sleeping place.)}

Ai cũng có thể bị thay.
\textit{(Anyone can be replaced.)}

Câu hỏi là em học gì trước sự thật đó.
\textit{(The question is what you learn from that fact.)}

\end{truyen}

\begin{baihoc}
Đừng sống như thể vị trí hôm nay là vĩnh viễn. Giữ mình sắc là cách duy nhất để không bị cũ quá nhanh.
\textit{(Do not live as if today’s place is permanent. Staying sharp is the only way not to become outdated too fast.)}
\end{baihoc}

\begin{ghinhoanh}
\textbf{Short English to remember:}

Do not live as if today’s place is permanent.

Staying sharp is the only way not to become outdated too fast.

\end{ghinhoanh}

\begin{tuhocnhanh}
\textbf{Cau hoi tu hoc / Self-study questions}

1. Van de chinh la gi? \textit{What is the main problem?}

2. Bai hoc thuc te la gi? \textit{What is the practical lesson?}

3. Mot cau tieng Anh em muon nho la cau nao? \textit{Which English sentence do you want to remember?}
\end{tuhocnhanh}
\ngancach

\section{Có Những Cánh Cửa Đóng Vì Em Chưa Đủ, Không Phải Vì Em Hỏng}
\begin{truyen}{Có Những Cánh Cửa Đóng Vì Em Chưa Đủ, Không Phải Vì Em Hỏng}{Some Doors Close Because You Are Not Ready Yet}
\chuhoa{M}{ột người trẻ bị từ chối vài lần rất dễ đi đến kết luận rằng mình không có tương lai.}
\textit{(A young person rejected several times can easily conclude that they have no future.)}

Đó là kết luận quá lớn từ dữ liệu còn quá ít.
\textit{(That is too large a conclusion from too little data.)}

Nhiều cánh cửa đóng không phải để sỉ nhục em.
\textit{(Many closed doors are not meant to humiliate you.)}

Nó chỉ phản ánh khoảng cách giữa em hiện tại và tiêu chuẩn của cơ hội đó.
\textit{(They only reflect the gap between your current self and that opportunity’s standard.)}

Thấy đúng khoảng cách thì còn học được.
\textit{(If you see the gap clearly, you can still learn.)}

Chỉ thấy mình hỏng thì thường chỉ muốn bỏ.
\textit{(If you see only yourself as broken, you will usually just want to quit.)}

\end{truyen}

\begin{baihoc}
Đừng biến “chưa đủ” thành “không xứng”. Hai câu đó khác nhau rất xa.
\textit{(Do not turn “not ready yet” into “not worthy.” Those are very different sentences.)}
\end{baihoc}

\begin{ghinhoanh}
\textbf{Short English to remember:}

Do not turn “not ready yet” into “not worthy.” Those are very different sentences.

\end{ghinhoanh}

\begin{tuhocnhanh}
\textbf{Cau hoi tu hoc / Self-study questions}

1. Van de chinh la gi? \textit{What is the main problem?}

2. Bai hoc thuc te la gi? \textit{What is the practical lesson?}

3. Mot cau tieng Anh em muon nho la cau nao? \textit{Which English sentence do you want to remember?}
\end{tuhocnhanh}
\ngancach

\section{Không Ai Nhớ Mãi Một Lần Em Trượt}
\begin{truyen}{Không Ai Nhớ Mãi Một Lần Em Trượt}{No One Remembers Your Failure Forever}
\chuhoa{N}{gười trẻ thường nghĩ lần thất bại của mình sẽ bám theo rất lâu trong mắt người khác.}
\textit{(Young people often think their failure will stay in other people’s eyes for a very long time.)}

Sự thật là ai cũng bận với cuộc đời của họ hơn em tưởng.
\textit{(The truth is that most people are more busy with their own lives than you imagine.)}

Điều bám lâu nhất không phải ký ức của người khác.
\textit{(What lasts longest is not other people’s memory.)}

Nó là cách em tự kể lại chuyện đó với chính mình.
\textit{(It is the story you keep telling yourself.)}

Nếu em không biến một cú trượt thành danh tính, nó sẽ sớm chỉ còn là một sự kiện.
\textit{(If you do not turn one failure into an identity, it soon becomes only an event.)}

\end{truyen}

\begin{baihoc}
Sau khi ngã, việc khó nhất không phải chịu mắt người đời. Nhiều khi là ngừng tự đóng đinh mình vào cú ngã đó.
\textit{(After a fall, the hardest task is often not facing the world. It is stopping yourself from nailing your whole identity to that fall.)}
\end{baihoc}

\begin{ghinhoanh}
\textbf{Short English to remember:}

After a fall, the hardest task is often not facing the world.

It is stopping yourself from nailing your whole identity to that fall.

\end{ghinhoanh}

\begin{tuhocnhanh}
\textbf{Cau hoi tu hoc / Self-study questions}

1. Van de chinh la gi? \textit{What is the main problem?}

2. Bai hoc thuc te la gi? \textit{What is the practical lesson?}

3. Mot cau tieng Anh em muon nho la cau nao? \textit{Which English sentence do you want to remember?}
\end{tuhocnhanh}
\ngancach

\section{Nhìn Người Khác Được Chọn}
\begin{truyen}{Nhìn Người Khác Được Chọn}{Watching Others Get Chosen}
\chuhoa{C}{ó người liên tục nhìn bạn bè hoặc đồng nghiệp được chọn còn mình thì bị bỏ qua.}
\textit{(Some people keep watching friends or coworkers get chosen while they are left behind.)}

Họ tưởng bị bỏ lại nhiều lần nghĩa là mình kém hẳn.
\textit{(They think repeated exclusion means they are truly inferior.)}

Nhưng có khi đó chỉ là giai đoạn người khác đang hợp hơn với một loại cơ hội nhất định.
\textit{(But sometimes it only means others fit a certain opportunity better for now.)}

Nếu không biến nó thành danh tính, em vẫn còn rất nhiều đường để đi tiếp.
\textit{(If you do not turn it into identity, many roads still remain open.)}

\end{truyen}

\begin{baihoc}
Bị chọn muộn không có nghĩa là không có lượt. Miễn là em vẫn giữ mình sắc.
\textit{(Being chosen later does not mean you have no turn. Keep yourself sharp.)}
\end{baihoc}

\begin{ghinhoanh}
\textbf{Short English to remember:}

Being chosen later does not mean you have no turn.

Keep yourself sharp.

\end{ghinhoanh}

\begin{tuhocnhanh}
\textbf{Cau hoi tu hoc / Self-study questions}

1. Van de chinh la gi? \textit{What is the main problem?}

2. Bai hoc thuc te la gi? \textit{What is the practical lesson?}

3. Mot cau tieng Anh em muon nho la cau nao? \textit{Which English sentence do you want to remember?}
\end{tuhocnhanh}
\ngancach

\section{Bị Góp Ý Trước Mặt Người Khác}
\begin{truyen}{Bị Góp Ý Trước Mặt Người Khác}{Being Corrected in Front of Others}
\chuhoa{C}{ó lần em bị chê ngay trong phòng họp và cảm thấy mặt mình nóng ran.}
\textit{(Once you are criticized in a meeting and feel your face burning.)}

Nhiều người tưởng khoảnh khắc đó đủ để làm mình mất giá luôn.
\textit{(Many think that moment permanently reduces their value.)}

Thực ra điều người ta nhớ lâu hơn thường là cách em đón nhận và sửa sau đó.
\textit{(In reality, what people often remember longer is how you receive and fix it afterward.)}

Một phản ứng bình tĩnh đôi khi cứu em nhiều hơn cả một phần giải thích dài.
\textit{(A calm response sometimes saves you more than a long explanation.)}

\end{truyen}

\begin{baihoc}
Bị chê trước mặt người khác đau, nhưng cách em đứng lại sau đó mới nói rõ em là ai.
\textit{(Public criticism hurts, but how you stand afterward says more about you.)}
\end{baihoc}

\begin{ghinhoanh}
\textbf{Short English to remember:}

Public criticism hurts, but how you stand afterward says more about you.

\end{ghinhoanh}

\begin{tuhocnhanh}
\textbf{Cau hoi tu hoc / Self-study questions}

1. Van de chinh la gi? \textit{What is the main problem?}

2. Bai hoc thuc te la gi? \textit{What is the practical lesson?}

3. Mot cau tieng Anh em muon nho la cau nao? \textit{Which English sentence do you want to remember?}
\end{tuhocnhanh}
\ngancach

\section{Thư Từ Chối Đầu Tiên}
\begin{truyen}{Thư Từ Chối Đầu Tiên}{The First Rejection Letter}
\chuhoa{M}{ột email từ chối đầu tiên có thể làm người trẻ ngồi lặng rất lâu.}
\textit{(The first rejection email can leave a young person silent for a long time.)}

Họ tưởng một câu không đủ nói rằng con đường đó không dành cho mình.
\textit{(They think one no is enough to prove the road is not theirs.)}

Nhưng rất nhiều cánh cửa chỉ đang nói rằng em chưa đúng lúc, chưa đúng cách, hoặc chưa đúng chỗ.
\textit{(But many doors are only saying you are not right yet, not right there, or not right now.)}

Bài học nằm ở việc đọc cú từ chối như dữ liệu chứ không như phán quyết toàn đời.
\textit{(The lesson is to read rejection as data, not as a lifelong verdict.)}

\end{truyen}

\begin{baihoc}
Một lá thư từ chối không viết xong giá trị của em. Nó chỉ đóng một cánh cửa hôm đó.
\textit{(One rejection letter does not write your full worth. It closes only one door that day.)}
\end{baihoc}

\begin{ghinhoanh}
\textbf{Short English to remember:}

One rejection letter does not write your full worth.

It closes only one door that day.

\end{ghinhoanh}

\begin{tuhocnhanh}
\textbf{Cau hoi tu hoc / Self-study questions}

1. Van de chinh la gi? \textit{What is the main problem?}

2. Bai hoc thuc te la gi? \textit{What is the practical lesson?}

3. Mot cau tieng Anh em muon nho la cau nao? \textit{Which English sentence do you want to remember?}
\end{tuhocnhanh}
\ngancach

\section{Đường Đi Thứ Hai}
\begin{truyen}{Đường Đi Thứ Hai}{The Second Choice Can Still Be a Road}
\chuhoa{C}{ó người không vào được nơi mình mơ nên miễn cưỡng bước vào một lựa chọn thứ hai.}
\textit{(Some people fail to enter their dream place and reluctantly step into a second choice.)}

Họ tưởng đường tốt nhất đã mất thì phần còn lại chỉ là chấp nhận thua.
\textit{(They think losing the best road means accepting defeat.)}

Nhưng nhiều đời người đi rất xa từ con đường không phải lựa chọn đầu tiên.
\textit{(But many lives travel far on roads that were not the first choice.)}

Quan trọng không chỉ là cửa em vào đầu tiên, mà là em đi tiếp trong cửa đó thế nào.
\textit{(What matters is not only the first door you enter, but how you walk inside it.)}

\end{truyen}

\begin{baihoc}
Lựa chọn thứ hai không tự động là đời hạng hai. Nó còn tùy em làm gì với nó.
\textit{(A second choice is not automatically a second-rate life. It depends on what you build there.)}
\end{baihoc}

\begin{ghinhoanh}
\textbf{Short English to remember:}

A second choice is not automatically a second-rate life.

It depends on what you build there.

\end{ghinhoanh}

\begin{tuhocnhanh}
\textbf{Cau hoi tu hoc / Self-study questions}

1. Van de chinh la gi? \textit{What is the main problem?}

2. Bai hoc thuc te la gi? \textit{What is the practical lesson?}

3. Mot cau tieng Anh em muon nho la cau nao? \textit{Which English sentence do you want to remember?}
\end{tuhocnhanh}
\ngancach

\section{Đứng Dậy Sau Khi Bị Thay}
\begin{truyen}{Đứng Dậy Sau Khi Bị Thay}{Standing Up After Being Replaced}
\chuhoa{B}{ị thay ở một vị trí làm nhiều người trẻ thấy như mình bị phủ nhận hoàn toàn.}
\textit{(Being replaced in a role makes many young people feel totally rejected.)}

Họ tưởng mất chỗ đó đồng nghĩa mình không còn giá trị.
\textit{(They think losing that place means they no longer matter.)}

Nhưng có lúc bị thay chỉ là lời nhắc rằng năng lực cần được mài lại, hoặc môi trường đó đã không còn hợp.
\textit{(Sometimes replacement only reminds you to sharpen your ability again, or it shows the place no longer fits.)}

Đứng dậy sau cú đó thường cho em một kiểu bản lĩnh sâu hơn cả lúc được giữ lại.
\textit{(Standing up after that blow often gives you a deeper strength than being kept would have.)}

\end{truyen}

\begin{baihoc}
Bị thay đau thật. Nhưng nếu học được, em sẽ không bị cũ đi cùng cú đau đó.
\textit{(Being replaced truly hurts. But if you learn from it, you will not grow old inside that pain.)}
\end{baihoc}

\begin{ghinhoanh}
\textbf{Short English to remember:}

Being replaced truly hurts.

But if you learn from it, you will not grow old inside that pain.

\end{ghinhoanh}

\begin{tuhocnhanh}
\textbf{Cau hoi tu hoc / Self-study questions}

1. Van de chinh la gi? \textit{What is the main problem?}

2. Bai hoc thuc te la gi? \textit{What is the practical lesson?}

3. Mot cau tieng Anh em muon nho la cau nao? \textit{Which English sentence do you want to remember?}
\end{tuhocnhanh}
